\documentclass{harryopo-report}

\title{基于深度学习的智慧交通流量预测系统设计报告}
\author{王明 计算机学院 2503010102}
\institute{示例大学}
\date{2026年08月30日}

\begin{document}

\maketitle

\tableofcontents

\chapter*{第一章 绪论}
\addcontentsline{toc}{chapter}{第一章 绪论}

\section*{1.1 研究背景与意义}
\addcontentsline{toc}{section}{1.1 研究背景与意义}

随着城市化进程的加速，交通拥堵已成为深圳等大型城市面临的严峻挑战。传统的交通流量预测方法——如时间序列分析和统计回归——难以捕捉复杂的时空依赖关系，在面对突发状况时预测精度大幅下降。

基于深度学习的智慧交通预测系统，能够通过路网拓扑建模和时序演化分析，提供更准确的预测结果，从而辅助交通管理部门进行科学决策。

\section*{1.2 国内外研究现状}
\addcontentsline{toc}{section}{1.2 国内外研究现状}

\subsection*{1.2.1 国外研究现状}
\addcontentsline{toc}{subsection}{1.2.1 国外研究现状}

国外学者较早开始研究交通流量预测问题。文献 [1] 提出基于 LSTM 的预测模型，在 METR-LA 数据集上取得良好效果。文献 [2] 进一步引入图卷积网络，建模路网空间结构。

\subsection*{1.2.2 国内研究现状}
\addcontentsline{toc}{subsection}{1.2.2 国内研究现状}

国内研究多聚焦于特定城市的实际场景。文献 [3] 针对北京市的交通特点提出了改进的 GRU 模型。文献 [4] 则研究了深圳的交通流特性，构建了融合多源数据的预测框架。

\section*{1.3 主要研究内容}
\addcontentsline{toc}{section}{1.3 主要研究内容}

本报告的主要研究内容包括：

\begin{enumerate}
  \item {\fzht 数据集构建}：采集深圳市福田区 6 个月的交通流量数据
  \item {\fzht 模型设计}：设计图卷积-时序混合网络架构
  \item {\fzht 系统实现}：开发完整的 B/S 架构 Web 系统
  \item {\fzht 部署验证}：在真实场景中验证系统有效性
\end{enumerate}

\section*{1.4 报告组织结构}
\addcontentsline{toc}{section}{1.4 报告组织结构}

本报告共分六章：

\begin{itemize}
  \item {\fzht 第一章} 绪论
  \item {\fzht 第二章} 相关技术与理论基础
  \item {\fzht 第三章} 需求分析与系统设计
  \item {\fzht 第四章} 模型设计与实现
  \item {\fzht 第五章} 系统测试与结果分析
  \item {\fzht 第六章} 总结与展望
\end{itemize}

\medskip\hrule\medskip

\chapter*{第二章 相关技术与理论基础}
\addcontentsline{toc}{chapter}{第二章 相关技术与理论基础}

\section*{2.1 深度学习基础}
\addcontentsline{toc}{section}{2.1 深度学习基础}

\subsection*{2.1.1 卷积神经网络（CNN）}
\addcontentsline{toc}{subsection}{2.1.1 卷积神经网络（CNN）}

卷积神经网络是处理网格结构数据（如图像）的经典架构。其核心思想是通过卷积核提取局部特征，再通过池化层降维。

\subsection*{2.1.2 循环神经网络（RNN）}
\addcontentsline{toc}{subsection}{2.1.2 循环神经网络（RNN）}

循环神经网络擅长处理序列数据。LSTM 通过门控机制解决了传统 RNN 的梯度消失问题。

\begin{lstlisting}[style=pystyle,caption={}]
import torch
import torch.nn as nn

class TrafficPredictor(nn.Module):
    def __init__(self, in_dim, hidden_dim, out_dim, num_layers=2):
        super().__init__()
        self.lstm = nn.LSTM(
            input_size=in_dim,
            hidden_size=hidden_dim,
            num_layers=num_layers,
            batch_first=True,
            dropout=0.2
        )
        self.fc = nn.Linear(hidden_dim, out_dim)

    def forward(self, x):
        # x: (batch, seq_len, in_dim)
        lstm_out, _ = self.lstm(x)
        return self.fc(lstm_out[:, -1, :])
\end{lstlisting}

\section*{2.2 图神经网络}
\addcontentsline{toc}{section}{2.2 图神经网络}

\subsection*{2.2.1 图卷积网络（GCN）}
\addcontentsline{toc}{subsection}{2.2.1 图卷积网络（GCN）}

图卷积网络将传统卷积操作推广到图结构数据，能够有效建模节点间的拓扑关系。

\begin{quote}
{\fzht 公式 2.1} GCN 层的传播规则：
\$$H^{(l+1)} = \sigma\left(\tilde{D}^{-1/2} \tilde{A} \tilde{D}^{-1/2} H^{(l)} W^{(l)}\right)$\$
\end{quote}

其中 $\tilde{A} = A + I$ 为加自环的邻接矩阵，$\tilde{D}$ 为对应的度矩阵。

\subsection*{2.2.2 图注意力网络（GAT）}
\addcontentsline{toc}{subsection}{2.2.2 图注意力网络（GAT）}

GAT 通过注意力机制自适应地学习邻居节点的权重，进一步提升性能。

\begin{table}[htbp]
  \centering
  \small
  \begin{tabularx}{\textwidth}{>{\raggedright\arraybackslash}X >{\raggedright\arraybackslash}X >{\raggedright\arraybackslash}X >{\raggedright\arraybackslash}X}
    \toprule
    {\fzht 模型} & {\fzht MAE} & {\fzht RMSE} & {\fzht MAPE (\%)} \\
    \midrule
    HA & 6.02 & 9.78 & 17.50 \\
    LSTM & 4.27 & 6.85 & 12.20 \\
    DCRNN & 3.05 & 5.12 & 8.90 \\
    Graph WaveNet & 2.69 & 4.65 & 7.60 \\
    {\fzht 本文模型} & {\fzht 2.31} & {\fzht 3.98} & {\fzht 6.45} \\
    \bottomrule
  \end{tabularx}
\end{table}

\section*{2.3 本章小结}
\addcontentsline{toc}{section}{2.3 本章小结}

本章系统梳理了深度学习与图神经网络的基础理论，为后续模型设计奠定了理论基础。

---

\chapter*{第三章 需求分析与系统设计}
\addcontentsline{toc}{chapter}{第三章 需求分析与系统设计}

\section*{3.1 功能性需求}
\addcontentsline{toc}{section}{3.1 功能性需求}

\subsection*{3.1.1 数据采集模块}
\addcontentsline{toc}{subsection}{3.1.1 数据采集模块}

\begin{itemize}
  \item 支持实时接入路口摄像头数据
  \item 支持 CSV/JSON 批量导入
  \item 数据清洗与异常值处理
\end{itemize}

\subsection*{3.1.2 预测服务模块}
\addcontentsline{toc}{subsection}{3.1.2 预测服务模块}

\begin{itemize}
  \item 支持未来 15/30/60 分钟多粒度预测
  \item 预测结果可视化展示
  \item 预测准确率实时监控
\end{itemize}

\section*{3.2 非功能性需求}
\addcontentsline{toc}{section}{3.2 非功能性需求}

\begin{itemize}
  \item {\fzht 性能}：单次预测响应时间 < 200ms
  \item {\fzht 可用性}：系统 7×24 小时稳定运行
  \item {\fzht 可扩展性}：支持横向扩展至 100+ 节点
\end{itemize}

\section*{3.3 系统架构设计}
\addcontentsline{toc}{section}{3.3 系统架构设计}

系统采用前后端分离的 B/S 架构：

\begin{lstlisting}[style=plainstyle,caption={}]
┌─────────────────────────────────────────────────┐
│               浏览器（Vue + ECharts）              │
└─────────────────┬───────────────────────────────┘
                  │ HTTPS / WebSocket
┌─────────────────▼───────────────────────────────┐
│          后端服务（FastAPI + Uvicorn）              │
│  ┌─────────┐ ┌──────────┐ ┌──────────────┐      │
│  │ 模型推理 │ │ 数据管理 │ │ 用户认证      │      │
│  └─────────┘ └──────────┘ └──────────────┘      │
└─────────────────┬───────────────────────────────┘
                  │ TCP / RPC
┌─────────────────▼───────────────────────────────┐
│         模型服务（PyTorch + TorchServe）            │
└─────────────────────────────────────────────────┘
\end{lstlisting}

\begin{quote}
{\fzht 架构特点}：前后端通过 RESTful API 解耦，模型服务可独立扩容，支持 GPU 集群推理。
\end{quote}

\medskip\hrule\medskip

\chapter*{第四章 模型设计与实现}
\addcontentsline{toc}{chapter}{第四章 模型设计与实现}

\section*{4.1 数据预处理}
\addcontentsline{toc}{section}{4.1 数据预处理}

\subsection*{4.1.1 数据集描述}
\addcontentsline{toc}{subsection}{4.1.1 数据集描述}

本报告使用深圳市福田区 2024 年 1 月至 6 月的交通流量数据，包含：

\begin{itemize}
  \item 156 个关键路口的实时车流量
  \item 5 分钟粒度采样，共约 50 万条记录
  \item 天气、节假日等外部特征
\end{itemize}

\subsection*{4.1.2 特征工程}
\addcontentsline{toc}{subsection}{4.1.2 特征工程}

\begin{itemize}
  \item 时序特征：小时、星期、是否节假日
  \item 空间特征：路口邻接关系、路段长度
  \item 外部特征：天气、温度、特殊事件
\end{itemize}

\section*{4.2 模型架构}
\addcontentsline{toc}{section}{4.2 模型架构}

\subsection*{4.2.1 整体架构}
\addcontentsline{toc}{subsection}{4.2.1 整体架构}

模型由{\fzht 空间建模}和{\fzht 时序建模}两部分组成：

\begin{quote}
{\fzht 空间建模}：使用 GAT 建模路口间的空间依赖关系

{\fzht 时序建模}：使用 GRU 建模时序演化规律

{\fzht 融合层}：拼接空间与时序特征，通过全连接层输出最终预测
\end{quote}

\subsection*{4.2.2 损失函数}
\addcontentsline{toc}{subsection}{4.2.2 损失函数}

本报告使用 Masked MAE 损失，对缺失值不参与梯度计算：

\begin{quote}
{\fzht 公式 4.1} Masked MAE 损失：
\$$L = \frac{\sum_{i} m_i \cdot |y_i - \hat{y}_i|}{\sum_{i} m_i + \epsilon}$\$
\end{quote}

其中 $m_i \in \{0, 1\}$ 为掩码标记，$\epsilon$ 为防止除零的小常数。

\section*{4.3 训练策略}
\addcontentsline{toc}{section}{4.3 训练策略}

\begin{itemize}
  \item {\fzht 优化器}：Adam（lr=1e-3, weight\_decay=1e-5）
  \item {\fzht 学习率调度}：ReduceLROnPlateau（factor=0.5, patience=5）
  \item {\fzht 批量大小}：64
  \item {\fzht 训练轮数}：100（early stopping，patience=10）
\end{itemize}

\medskip\hrule\medskip

\chapter*{第五章 系统测试与结果分析}
\addcontentsline{toc}{chapter}{第五章 系统测试与结果分析}

\section*{5.1 测试环境}
\addcontentsline{toc}{section}{5.1 测试环境}

\begin{table}[htbp]
  \centering
  \small
  \begin{tabularx}{\textwidth}{>{\raggedright\arraybackslash}X >{\raggedright\arraybackslash}X}
    \toprule
    {\fzht 组件} & {\fzht 配置} \\
    \midrule
    CPU & Intel Xeon Gold 6248 @ 2.5GHz × 2 \\
    GPU & NVIDIA Tesla V100 32GB \\
    内存 & 128GB DDR4 \\
    操作系统 & Ubuntu 20.04 LTS \\
    深度学习框架 & PyTorch 2.0 + CUDA 11.8 \\
    \bottomrule
  \end{tabularx}
\end{table}

\section*{5.2 性能测试}
\addcontentsline{toc}{section}{5.2 性能测试}

\subsection*{5.2.1 预测精度}
\addcontentsline{toc}{subsection}{5.2.1 预测精度}

\begin{table}[htbp]
  \centering
  \small
  \begin{tabularx}{\textwidth}{>{\raggedright\arraybackslash}X >{\raggedright\arraybackslash}X >{\raggedright\arraybackslash}X >{\raggedright\arraybackslash}X}
    \toprule
    {\fzht 预测时长} & {\fzht MAE} & {\fzht RMSE} & {\fzht MAPE (\%)} \\
    \midrule
    15 分钟 & 1.92 & 3.21 & 5.12 \\
    30 分钟 & 2.31 & 3.98 & 6.45 \\
    60 分钟 & 2.87 & 4.76 & 8.13 \\
    \bottomrule
  \end{tabularx}
\end{table}

\subsection*{5.2.2 系统响应时间}
\addcontentsline{toc}{subsection}{5.2.2 系统响应时间}

\begin{table}[htbp]
  \centering
  \small
  \begin{tabularx}{\textwidth}{>{\raggedright\arraybackslash}X >{\raggedright\arraybackslash}X >{\raggedright\arraybackslash}X}
    \toprule
    {\fzht 并发数} & {\fzht 平均响应 (ms)} & {\fzht P99 响应 (ms)} \\
    \midrule
    10 & 45 & 89 \\
    50 & 78 & 156 \\
    100 & 124 & 245 \\
    200 & 198 & 387 \\
    \bottomrule
  \end{tabularx}
\end{table}

\section*{5.3 案例分析}
\addcontentsline{toc}{section}{5.3 案例分析}

以福田区深南大道-华富路口为例：

\begin{itemize}
  \item {\fzht 工作日早高峰}（8:00-9:00）：预测车流量 1,250 辆/小时，实际 1,287 辆/小时（误差 2.9\%）
  \item {\fzht 周末晚高峰}（18:00-19:00）：预测 980 辆/小时，实际 952 辆/小时（误差 2.9\%）
  \item {\fzht 特殊事件}（2024 年跨年夜）：预测 2,150 辆/小时，实际 2,398 辆/小时（误差 10.4\%，仍优于基线）
\end{itemize}

\medskip\hrule\medskip

\chapter*{第六章 总结与展望}
\addcontentsline{toc}{chapter}{第六章 总结与展望}

\section*{6.1 工作总结}
\addcontentsline{toc}{section}{6.1 工作总结}

本报告完成的主要工作包括：

\begin{enumerate}
  \item 系统梳理了智慧交通预测领域的国内外研究现状
  \item 构建了基于图卷积-时序网络的混合预测模型
  \item 设计并实现了完整的 B/S 架构 Web 预测系统
  \item 在深圳市福田区真实场景中验证了系统有效性
\end{enumerate}

\section*{6.2 主要创新点}
\addcontentsline{toc}{section}{6.2 主要创新点}

\begin{itemize}
  \item {\fzht 多尺度时空融合}：同时建模短期突变与长期趋势
  \item {\fzht 外部特征自适应融合}：自动学习天气、节假日等外部因素的影响
  \item {\fzht 边缘-云端协同部署}：在保证精度的同时提升响应速度
\end{itemize}

\section*{6.3 不足与展望}
\addcontentsline{toc}{section}{6.3 不足与展望}

本系统仍存在以下不足：

\begin{enumerate}
  \item 模型可解释性有待加强
  \item 极端天气下的预测精度需进一步提升
  \item 暂未支持跨城市迁移
\end{enumerate}

未来将探索：

\begin{itemize}
  \item 引入因果推断提升可解释性
  \item 研究联邦学习支持跨城市协同训练
  \item 集成大语言模型实现自然语言交互式查询
\end{itemize}

\medskip\hrule\medskip

\begin{thebibliography}{99}
\bibitem{ref1} Li Y, Yu R, Shahabi C, et al. Diffusion convolutional recurrent neural network: Data-driven traffic forecasting[C]. ICLR, 2018.
\bibitem{ref2} Zhao L, Song Y, Zhang C, et al. T-GCN: A temporal graph convolutional network for traffic prediction[J]. IEEE Transactions on Intelligent Transportation Systems, 2019, 21(9): 3848-3858.
\bibitem{ref3} 王伟, 李强. 基于改进 GRU 的城市交通流量预测方法[J]. 计算机应用研究, 2022, 39(5): 1452-1456.
\bibitem{ref4} 陈晓, 刘洋. 融合多源数据的深圳交通流预测研究[J]. 智能系统学报, 2023, 18(3): 567-575.
\end{thebibliography}

\end{document}